% !PRINTED 260
\weberpage{260}
%<<CONT
and $\delta$ and $\beta$ can be so determined that

\weq{7}{\delta(\alpha+\beta) - \beta(\gamma+\delta) = \alpha\delta - \beta\gamma = 1.}

From the equations (6) for $a$ and $c-\partial$ there is then found

\weq{8}{c'+a' = a\delta - (c-\partial)\beta.}

Since $\alpha+\beta$ and $\gamma+\delta$ are both odd, $\delta$ and $\beta$
cannot, by (7), both be odd; consequently $c'$ turns out even by (8).
If one replaces $\delta$, $\beta$, which by (7) is permitted, by
$\delta+h(\gamma+\delta)$, $\beta+h(\alpha+\beta)$, for an arbitrary $h$, then
$c'$ changes by $2ha'$ according to (6) and (8), and $h$ can be so disposed
of that $c'$ becomes divisible by 16 or 48.

Then, by (6), $\alpha+\beta+\gamma-\delta$ is even, and therefore
formula (12), \S\,40 is to be applied. In it the following is now to be taken into account


\[
\begin{aligned}
a'(\alpha-\beta) &\equiv a, & \partial(\alpha-\beta) &\equiv \partial' \\
\partial'(\alpha+\beta) &\equiv a, & a'(\gamma-\delta) &\equiv -\partial
\end{aligned} \quad (\mathrm{mod}\ 16),
\]

from which it follows:

\[
n(\alpha-\beta)(\alpha+\beta+\gamma-\delta) \equiv a^2 - \partial'^2 \ (\mathrm{mod}\ 16),
\]

hence

\[
\left(\frac{2}{\alpha-\beta}\right) e^{-\frac{3\pi i}{8}(\alpha-\beta)(\alpha+\beta+\gamma-\delta)} = \left(\frac{2}{a}\right)\left(\frac{2}{a'}\right)\left(\frac{2}{a}\right)\left(\frac{2}{\partial'}\right) = \left(\frac{2}{n}\right)
\]

and

\[
\varrho = e^{-\frac{2\pi i}{3}\left[\alpha(\gamma-\beta)-(\alpha^2-1)\beta\delta\right]}.
\]

If $n$ is not divisible by 3, then

\[
\begin{aligned}
2\alpha &\equiv a(a'+\partial'), & 2\delta &\equiv \partial(a'+\partial'), \\
2\beta &\equiv a(\partial'-a'), & 2\gamma &\equiv \partial(\partial'-a')
\end{aligned} \quad (\mathrm{mod}\ 3),
\]

hence either $\alpha$ and $\delta$ or $\beta$ and $\gamma$ are divisible by 3, and therefore
$\varrho = 1$.

From this there result the following interchanges belonging together
 (where the change of $a, c, \partial$ into $a', c', \partial'$ is determined by (6)):

\weq{9}{\begin{array}{ccc}
\omega, & u, & v \\[4pt]
\dfrac{\omega-1}{\omega+1}, & \dfrac{\sqrt{2}}{u}, & \varrho\left(\dfrac{2}{n}\right)\dfrac{\sqrt{2}}{v}.
\end{array}}

We apply the interchanges (4), (5), (9) to the following
functions:
%CONT>>
